\chapter*{ВВЕДЕНИЕ}
\addcontentsline{toc}{chapter}{ВВЕДЕНИЕ}

Интерполяция видеокадров (Video Frame Interpolation, VFI)~---~это синтез промежуточных кадров между двумя соседними кадрами видеопоследовательности. Технология применяется для повышения частоты кадров, формирования эффекта замедленного воспроизведения и сокращения объёма данных при сжатии видео.

Нейросетевые методы интерполяции обеспечивают высокое качество синтеза, однако точное восстановление промежуточного кадра остаётся затруднительным при нелинейном движении объектов и наличии окклюзий. Современные архитектуры, в частности IFRNet, достигают высоких показателей качества, но характеризуются значительной вычислительной сложностью. Обработка изображений в цветовом пространстве RGB сопряжена с высокими затратами ресурсов, что ограничивает применение таких моделей в системах реального времени.

В связи с этим разработка методов оптимизации, снижающих вычислительную нагрузку без существенной потери визуального качества, в том числе с учётом особенностей зрительного восприятия человека, представляет собой актуальную задачу.

Целью работы является сравнительный анализ методов оптимизации архитектуры IFRNet, направленных на повышение вычислительной эффективности.

Для достижения поставленной цели необходимо решить следующие задачи:
\begin{enumerate}
	\item провести анализ архитектуры IFRNet и обосновать её выбор в качестве объекта оптимизации среди современных методов интерполяции видеокадров;
	\item реализовать методы оптимизации: равномерное сокращение числа каналов (RGB/2) и асимметричную обработку в цветовом пространстве YUV420;
	\item обучить модели на наборе данных Vimeo90K с использованием функций потерь L1, Ternary и Distill;
	\item провести сравнительный анализ методов по показателям вычислительной эффективности (FLOPs, число параметров, FPS) и качества синтеза (PSNR, SSIM, LPIPS, DISTS).
\end{enumerate}
